
\section{A General Continuum Solution}

The treatment follows the program of Ie\c{s}an, whereby the Saint Venant class is decomposed into two problems: that of torsion-bending-elongation, and flexure.


Consider a prismatic, homogeneous, and isotropic linear elastic cylinder, $\Body$. 
The body occupies the region $\Body = \Section \times (0, L)$, where $\Section$ is a region in the $\Real^2$ plane spanned by the orthonormal basis $\{\FrameBasis_y, \FrameBasis_z\}$ and $L$ is the cylinder's length along the $\FrameBasis$-axis. 
The vectors \(\{\FrameBasis,\,\FrameBasis_y,\,\FrameBasis_z\}\) form an orthonormal basis in 3D.
$\Section$ may be multiply connected. 

The basis vector along the axis of the beam is $\FrameBasis$, and the section plane is spanned by an orthogonal basis $\{\FrameBasis_y,\,\FrameBasis_z\}$. 
Coordinates in the section are $\bm{r} = y\,\FrameBasis_y + z\,\FrameBasis_z$. 
The 3D identity is 
\[\mathbf{1} = \FrameBasis \otimes \FrameBasis + \mathbf{1}_\Section,
\quad\text{ with }\quad 
\mathbf{1}_\Section = \FrameBasis_y \otimes \FrameBasis_y + \FrameBasis_z \otimes \FrameBasis_z
\]
The distinction between the 2D and 3D entity will be clear from context so distinct symbols for $\bm{r}$ need not be introduced. Similarly allow $\FrameBasis$ to be regarded as a 2D vector when appropriate to form terms like $\FrameBasis\times\bm{r}$. 

The lateral surface is $\BodyLateralSurface = \SectionBoundary \times (0, L)$, where $\SectionBoundary$ is the boundary of the section $\Section$. 


\begin{remark}
The following developments are made without appealing to a coordinate system (we only assume right-handedness), and thus at no commitments are made to a particular coordinate convention. However, such conventions can aid the interpretation of equations and are required for implementation. 
Two common conventions for assigning coordinates \(x,y,z\) are:
\begin{itemize}
    \item Take the first (\(x\)) coordinate to coincide with the reference axis \(\FrameBasis\) of the rod. This furnishes components on the basis \(\{\bm{\partial}_x,\bm{\partial}_y,\bm{\partial}_z\}\)
    \begin{equation}
    \FrameBasis \equiv \bm{\partial}_x = \begin{pmatrix}
        1 \\ 0 \\ 0
    \end{pmatrix},
    \quad
    \PlanePosition = \begin{pmatrix}
    0 \\ y \\ z
    \end{pmatrix},
    \quad 
    \bm{r}^{\times} =\left[\begin{array}{ccc}
    0 & -z & y \\
    z & 0 &  0 \\
    -y & 0 & 0
    \end{array}\right]
    \quad\text{and}\quad 
    \FrameBasis^{\times} = \begin{bmatrix}
        0 & 0 & 0 \\
        0 & 0 & -1 \\
        0 & 1 & 0
    \end{bmatrix}
    \label{eq:define-basis-1}
    \end{equation}
    \item Take the last (\(z\)) coordinate to coincide with \(\FrameBasis\). This furnishes components on the basis \(\{\bm{\partial}_x,\bm{\partial}_y,\bm{\partial}_z\}\)
    
    \[
    \FrameBasis \equiv \bm{\partial}_z = \begin{pmatrix}
        0 \\ 0 \\ 1
    \end{pmatrix},
    \quad
    \PlanePosition = \begin{pmatrix}
    x \\ y \\ 0
    \end{pmatrix}
    \qquad\text{and}\qquad 
    \FrameBasis^{\times} = \begin{bmatrix}
        0 & -1 & 0 \\
        1 &  0 & 0 \\
        0 & 0 & 0
    \end{bmatrix}
    \]
\end{itemize}
\end{remark}

\subsection{Saint Venant Problems}

The standard boundary value problem of elastostatics seeks a displacement field $\disp$ such that:
\begin{equation}
\left\{
\begin{aligned}
\operatorname{div}\bm{T}&=\mathbf{0},\\
\bm{T}&=\lambda(\operatorname{tr}\bm{\varepsilon})\mathbf{1}+2G\,\bm{\varepsilon}, \\
\bm{\varepsilon}&=\tfrac12(\nabla\bm{u}+\nabla\bm{u}^{\!\top}).
\end{aligned}\right.
\label{eq:3D-LE}
\end{equation}
subject to traction or displacement boundary conditions.

Saint~Venant's problem, in its relaxed form, replaces the pointwise prescription of tractions on the ends, $\Section_0$ and $\Section_L$, with stipulations on the net stress resultants (force $\Force$ and moment $\Moment$) \citep{iesan1986saintvenants}. 
A key observation is that this relaxed problem does not possess a unique solution.

Ie\c{s}an's analysis provides a rational derivation of the canonical Saint-Venant solutions by decomposing the relaxed problem. The problem is split into $(P_1)$, covering extension, bending, and torsion (where the resultant shear force is zero), and $(P_2)$, covering flexure (where the resultant axial force and all moments are zero at the loaded end).

In both problems the Saint Venant stress hypothesis \eqref{eq:saint-venant-stress} is recovered:
\begin{equation*}
\StressTensor = \AxialStress \FrameBasis \otimes\FrameBasis + \ShearStress \otimes \FrameBasis
+\FrameBasis \otimes \ShearStress
\qquad\text{with}\qquad
\begin{aligned}
\AxialStress &\triangleq \FrameBasis\cdot \StressTensor\,\FrameBasis \\
\bm{\tau} &\triangleq \PlaneBasis_{\beta}\otimes \PlaneBasis_{\beta} \StressTensor\,\FrameBasis
\end{aligned}
% \label{eq:saint-venant-stress}
\end{equation*}
where \(\AxialStress\) is the scalar axial stress and \(\ShearStress \in \mathbb{R}^2\) the vector of shear stresses. 

Saint-Venant's problem consists in the determination of an equilibrium displacement field \(\bm{u}\) on \(\Body\), subject to the requirements

\begin{equation}
\bm{\Traction}(\bm{u})=\mathbf{0} \quad \text { on } \BodyLateralSurface, 
\quad \bm{\Traction}(\bm{u})=\hat{\bm{\Traction}}^{(\alpha)} \quad \text{ on } \Section_\alpha \quad(\alpha=1,2),
\label{eq:iesan-2-1}
\end{equation}

where \(\hat{\bm{\Traction}}^{(\alpha)}\) is a vector-valued function
preassigned on \(\Section_\alpha\). 
Necessary conditions for the existence of a solution of this problem are given by

\[
\int_{\dot{\Sigma}_1} \hat{\bm{\Traction}}^{(1)} d a+\int_{\dot{\Sigma}_2} \hat{\bm{\Traction}}^{(2)} d a=\mathbf{0}, 
\qquad
\int_{\dot{\Sigma}_1} \bm{r} \times \hat{\bm{\Traction}}^{(1)} d a+\int_{\dot{\Sigma}_2} \bm{r} \times \hat{\bm{\Traction}}^{(2)} d a=\mathbf{0},
\]

In the relaxed formulation of Saint-Venant's problem the conditions \eqref{eq:iesan-2-1} are replaced by

\begin{equation*}
\bm{\Traction}(\bm{u})=\mathbf{0} \quad \text { on } \BodyLateralSurface, 
\quad \mathscr{F}(\bm{u})=\ShearForce, \quad \mathscr{M}(\bm{u})=\bm{M},
% \label{eq:2-2}
\end{equation*}
%
where \(\ShearForce\) and \(\bm{M}\) are prescribed vectors representing the
resultant force and the resultant moment about \(O\) of the tractions
acting on \(\Section_1\). Accordingly, \(\mathscr{F}(\cdot)\) and
\(\mathscr{M}(\cdot)\) are the vector-valued linear functionals defined by 
\begin{equation}
\mathscr{F}(\bm{u})=\int_{\dot{\Sigma}_1} \bm{\Traction}(\bm{u}) d a, 
\quad 
{\mathscr{ M }}(\bm{u})=\int_{\dot{\Sigma}_1} \bm{r} \times \bm{\Traction}(\bm{u}) d a
\label{eq:def-resultants}
\end{equation}

On the face \(\Section_1\) the unit normal is \(\bm{n} = -\FrameBasis\), \eqref{eq:def-resultants} appears as

\begin{equation*}
\VectorNotation{
\begin{aligned}
\mathscr{F}(\bm{u})&=-\int_{\Section_{\mathbf{1}}} \bm{S}(\bm{u})\FrameBasis d a \\
\mathscr{M}_{\Section}(\bm{u})&=-\int_{\Section_1} \sigma(\bm{u})\, (\FrameBasis\times \bm{r})\, da . \\
\FrameBasis\cdot \mathscr{M}(\bm{u})&=-\int_{\Section_1} (\FrameBasis\times \bm{r})\cdot \bm S_{3\kappa}(\bm{u})\, da ,
\end{aligned}}
% \label{eq:2-4}
\end{equation*}

\subsection{Problem 1 (extension-bending-torsion)}
\subsubsection{Statement}

Find the elastostatic state $(\disp, \StressTensor)$ for the class \(P_1\) of extension-bending-torsion. 
The problem is uniquely defined by the following conditions.

\begin{enumerate}
    \item \textbf{Equilibrium:} The stress tensor $\StressTensor$ is divergence-free in the body $\Body$:
    \begin{equation}
    \left\{
        \begin{aligned}
        \operatorname{div} \StressTensor &= \mathbf{0} & \quad\text{in } \Body \\
        \StressTensor \BodyLateralNormal&= \mathbf{0}& \quad \text{on } \BodyLateralSurface \\
        \StressTensor &= \lambda \operatorname{tr}(\StrainTensor) \One + 2G \StrainTensor
        \end{aligned}
    \right.
    \end{equation}

    \item \textbf{Saint Venant Class:} 
    \[
    F_\alpha=0
    \]

    \item \textbf{Ie\c{s}an's Postulate}
    The axial derivative of the displacement is rigid on each section:
    \begin{equation}
    \bm{u}'= a_{\varepsilon} \FrameBasis + (a_{\vartheta} \FrameBasis -\FrameBasis\times\bm{a}_{\Sigma}) \times \bm{r},
    % \label{eq:Iesan-postulate-1}
    \end{equation}
    % or equivalently:
    % \begin{equation}
    % \bm{u}'= \bm{\alpha}+\bm{\beta} \times \bm{x},
    % \end{equation}
    Equivalently, stresses depend at most linearly on $\reflenx$ and the torsional amplitude vanishes.
\end{enumerate}

\subsubsection{Warping}
\begin{equation}
\boxed{
\begin{array}{rll}
\operatorname{div}_G\left(\left.\nabla \WarpTwistIesan+\FrameBasis \times\PlanePosition\right.\right) &=0 & \text { in } \Section, \\
\partial_{\bm{\nu}} \WarpTwistIesan &=-\FrameBasis\times\PlanePosition \cdot \bm{\nu} & \text { on } \SectionBoundary, \\
\int_{\Section} \WarpTwistIesan\, \volE &= \AvgSV
\end{array}
}
% \label{eq:bvp-warp-twist-iesan}
\end{equation}

\subsubsection{Displacement}
\[
\bm{w}=\WarpTwistIesan\FrameBasis  \, a_{\vartheta}-\nu\bm{r} \, a_{\varepsilon}
-\nu\operatorname{dev}\bm{r}\otimes\bm{r}\;\IesanLoadM
\]
From \eqref{eq:2-13}:
\begin{equation}
\begin{aligned}
    \hat{\bm{u}}\{\bm{a}\} &= \left(-\tfrac{1}{2}x^2 \mathbf{1} - \nu\operatorname{dev}\bm{r}\otimes\bm{r} + x\, \FrameBasis\otimes\bm{r}\right)\IesanLoadM
    \\ 
    &+ (x\FrameBasis - \nu \bm{r})a_{\varepsilon} \\
    &+(x\; \FrameBasis\times \bm{r} + \WarpTwistIesan\, \FrameBasis)a_{\vartheta}
\end{aligned}
% \label{eq:iesan-displacement-1}
\end{equation}

\subsubsection{Strain}


% \begin{Quote}

% r^2=\bm{r}\cdot\bm{r}, \bm{r}^\perp:=\FrameBasis\times\bm{r}, and
% \mathbf{1}_\Section is the in–plane identity (on \mathrm{span}\{\FrameBasis_y,\FrameBasis_z\}).
% Gradients split as \nabla=\FrameBasis\otimes\partial_x+\nabla_\Section.
% For a section–scalar \psi(\bm{r}), \nabla (\psi\,\FrameBasis)=\FrameBasis\otimes\nabla_\Section\psi.
% For a section–vector \bm v(\bm{r}), \gradS (\FrameBasis\,(\bm v\cdot\cdot))=\FrameBasis\otimes\nabla_\Section(\bm v\cdot\cdot).
% The infinitesimal strain is \bm{\varepsilon}=\sym(\nabla\hat{\bm{u}}):=\tfrac12(\nabla\hat{\bm{u}}+(\nabla\hat{\bm{u}})^{\mathrm t}).

\textbf{Displacement gradient.}
Write \(\hat{\bm{u}}=\hat{\bm{u}}^{(a)}+\hat{\bm{u}}^{(\varepsilon)}+\hat{\bm{u}}^{(a_\vartheta)}\) from \eqref{eq:iesan-displacement-1}.

\medskip
\noindent\emph{(i) The \(\IesanLoadM\)-term.}
\[
\hat{\bm{u}}^{(a)}
=\bigl(-\tfrac12 x^2\mathbf{1}_\Section-\nu\,\operatorname{dev}(\bm{r}\otimes\bm{r})+x\,\FrameBasis\otimes\bm{r}\bigr)\IesanLoadM.
\]
Using the identities furnishes
\[
\begin{aligned}
\nabla\hat{\bm{u}}^{(a)}
&= -x\,\FrameBasis\otimes\IesanLoadM
\;-\nu\,(\bm{r}\cdot\IesanLoadM)\,\mathbf{1}_\Section
\;+\ (\FrameBasis\otimes\FrameBasis)\,(\bm{r}\cdot\IesanLoadM)
\;+\ x\,\FrameBasis\otimes\IesanLoadM\\
&=(\bm{r}\cdot\IesanLoadM)\,\big(\FrameBasis\otimes\FrameBasis-\nu\,\mathbf{1}_\Section\big).
\end{aligned}
\]
Hence
\begin{equation}
\StrainTensor^{(a)}
=(\bm{r}\cdot\IesanLoadM)\,\big(\FrameBasis\otimes\FrameBasis-\nu\,\mathbf{1}_\Section\big).
\label{eq:iesan-strain-tensor-bending}
\end{equation}

\medskip
\noindent\emph{(ii) The \(a_{\varepsilon}\)-block.}
\[
\hat{\bm{u}}^{(\varepsilon)}=(x\,\FrameBasis-\nu\,\bm{r})\,a_\varepsilon
\;\;\Longrightarrow\;\;
\nabla\hat{\bm{u}}^{(\varepsilon)}=a_\varepsilon\big(\FrameBasis\otimes\FrameBasis-\nu\,\mathbf{1}_\Section\big).
\]
Thus
\begin{equation}
\bm{E}^{(\varepsilon)}
=a_\varepsilon\big(\FrameBasis\otimes\FrameBasis-\nu\,\mathbf{1}_\Section\big).
\label{eq:iesan-strain-tensor-elongation}
\end{equation}

\medskip
\noindent\emph{(iii)}
The \(a_{\vartheta}\)-term 
\(
\hat{\bm{u}}^{(\vartheta,c_\vartheta)}
=\big(x\,\FrameBasis\times\bm {r}+\varphi\,\FrameBasis\big)a_\vartheta
\)
of \eqref{eq:iesan-displacement-1} furnishes
\[
\begin{aligned}
\nabla\hat{\bm{u}}^{(\vartheta,c_\vartheta)}
&= a_\vartheta\Big(\FrameBasis\otimes(\FrameBasis\times\bm{r})+x\,\bm\epsilon_\Section+\FrameBasis\otimes\nabla \WarpTwistIesan\Big)
\end{aligned}
\]
Since \(\sym(\bm\epsilon_\Section)=\mathbf{0}\),
\begin{equation}
\StrainTensor^{(a_\vartheta)}
=\tfrac12\Big(
\FrameBasis\otimes\bm{\varepsilon}_{\vartheta}
+\bm{\varepsilon}_{\vartheta}\otimes\FrameBasis
\Big),
\quad
\bm{\varepsilon}_{\vartheta}:= a_\vartheta\,(\FrameBasis\times\bm{r})
+a_\vartheta\,\nabla \WarpTwistIesan.
\label{eq:iesan-strain-tensor-torsion-1}
\end{equation}

Summing \eqref{eq:iesan-strain-tensor-elongation}, \eqref{eq:iesan-strain-tensor-bending}, \eqref{eq:iesan-strain-tensor-torsion-1}:
\begin{equation}
\begin{aligned}
\bm{E}^{(a)}
&= (a_{\varepsilon} + \bm{r}\cdot\IesanLoadM)\,\big(\FrameBasis\otimes\FrameBasis-\nu\,\mathbf{1}_\Section\big)
+\ \tfrac12\Big(\FrameBasis\otimes\bm{\varepsilon}_{\vartheta}+\bm{\varepsilon}_{\vartheta}\otimes\FrameBasis\Big) \\[2pt]
\end{aligned}
\label{eq:iesan-strain-tensor-1}
\end{equation}

% \end{Quote}


\subsubsection{Stress}

\begin{Quote}
% \mathbf{1}=\FrameBasis\otimes\FrameBasis+\mathbf{1}_\Section, 
% \mathbf{1}_\Section=\FrameBasis_y\otimes\FrameBasis_y+\FrameBasis_z\otimes\FrameBasis_z.
% Linear elastic isotropy: \bm{\sigma}=2G\,\bm{\varepsilon}+\lambda\,\operatorname{tr}(\bm{\varepsilon})\,\mathbf{1}
% with G the shear modulus and \lambda the Lamé parameter
% \big(\lambda=\tfrac{2G\nu}{1-2\nu}=\tfrac{E\nu}{(1+\nu)(1-2\nu)}\big).

\medskip
\noindent\emph{(i) The \(\IesanLoadM\)-block.}
From \eqref{eq:iesan-strain-tensor-bending} with \(\operatorname{tr}\StrainTensor^{(a)}=(1-2\nu)(\bm{r}\cdot\IesanLoadM)\),
\[
\begin{aligned}
\StressTensor^{(a)}
&=2G\Big(\FrameBasis\otimes\FrameBasis-\nu\,\mathbf{1}_\Section\Big)(\bm{r}\cdot\IesanLoadM)
+\lambda(1-2\nu)(\bm{r}\cdot\IesanLoadM)\,\mathbf{1}\\
&=2G(1+\nu)\,(\bm{r}\cdot\IesanLoadM)\,\FrameBasis\otimes\FrameBasis.
\end{aligned}
\]
Thus, under pure bending the stress is purely uniaxial.

\medskip
\noindent\emph{(ii) The \(a_\varepsilon\)-block.}
From \eqref{eq:iesan-strain-tensor-elongation} with \(\operatorname{tr}\StrainTensor^{(\varepsilon)}=(1-2\nu)a_\varepsilon\),
\[
\StressTensor^{(\varepsilon)}
=2G(1+\nu)\,a_\varepsilon\,\FrameBasis\otimes\FrameBasis.
\]

\medskip
\noindent\emph{(iii) The \((a_\vartheta,c_\vartheta)\)-block.}
From \eqref{eq:iesan-strain-tensor-torsion-1} 
with \(\bm{\varepsilon}_{\vartheta}\) in–section and \(\operatorname{tr}\bm{E}^{(a_\vartheta)}=0\):
\[
\StressTensor^{(a_\vartheta)}
=2G\,\bm{E}^{(a_\vartheta)}
=G\Big(\FrameBasis\otimes\bm{\varepsilon}_{\vartheta}+\bm{\varepsilon}_{\vartheta}\otimes\FrameBasis\Big).
\]
\end{Quote}

\begin{Quote}
\begin{equation*}
\begin{aligned}
\sigma (\bm{v})&=(\lambda+2 G)\left(\bm{r}\cdot\IesanLoadM +a_{\varepsilon}\right) + \lambda \operatorname{div}_{\Section}\bm{w},\\
&=(\lambda+2 G)\left(\bm{r}\cdot\IesanLoadM +a_{\varepsilon}\right) - 2\lambda \nu\, (\bm{r}\cdot\IesanLoadM + a_{\varepsilon}),\\
&=E\big(a_{\varepsilon} + \bm{r}\cdot\IesanLoadM\big)
\end{aligned}
\end{equation*}
\end{Quote}
\begin{equation}
\begin{aligned}
\bm{\tau}(\bm{v})&=G\left(\nabla \WarpTwistIesan- \FrameBasis\times\bm{r}\right)a_{\vartheta}, \\ 
\sigma (\bm{v})
&=E\big(\bm{r}\cdot\IesanLoadM +a_{\varepsilon}\big)
\end{aligned}
\label{eq:iesan-stress-1}
\end{equation}

\subsubsection{Resultants}
%
The conditions on the terminal cross-section \(\Section_1\) furnish the
following system for the unknown constants

\begin{equation}
{
\begin{aligned}
E\left(\IntRoR \IesanLoadM + A \CentroidLocation a_{\varepsilon}\right)&=\FrameBasis\times \bm{M}, \quad\text{(sign?)}\\
E A\left(\IesanLoadM\cdot\CentroidLocation +a_{\varepsilon}\right)&=-N, \\
G \Jsv a_{\vartheta}&=-T,
\end{aligned}}
\label{eq:2-15}
\end{equation}

where the torsion constant \(\Jsv\) is defined by
\begin{equation}
\Jsv 
\triangleq \int [\FrameBasis,\, \bm{r},\, \nabla\WarpTwistIesan + \FrameBasis\times\bm{r}] \volG 
\equiv\int_{\Section} \Big(r^2 + (\FrameBasis\times \bm{r})\cdot \nabla \varphi\Big)\, \volG
\label{eq:def-j-origin}
\end{equation}


\clearpage
\subsection{Problem 2}
\subsubsection{Synopsis}
A central theorem states that if $\disp$ is a solution to the flexure problem $(\mathrm{P}_2)$ for a resultant shear force $\Force$, then its axial derivative, $\bm{u}'$, is a solution to the bending problem $(\mathrm{P}_1)$ corresponding to a resultant moment $\Moment = \FrameBasis \times \Force$ \citep{iesan1986saintvenants}. 
The canonical Saint-Venant solution for bending, $\bm{v}$, is itself uniquely characterized by the property that its second axial derivative, $\bm{v}''$, is zero (i.e., $\bm{v}'$ is a rigid displacement field). 
The flexure solution $\disp$ is then constructed by integrating this bending solution $\bm{v}$.

\subsubsection{Statement}
We seek the elastostatic state $(\disp, \StressTensor)$ for the torsion-free flexure of a cylinder $\Body = \Section \times (0, L)$, subject to a transverse force at one end. The problem is defined by the following conditions.

\begin{enumerate}
    \item \textbf{Equilibrium:} The stress tensor $\StressTensor$ is divergence-free in the body $\Body$:
    \begin{align*}
        \operatorname{div} \StressTensor &= \mathbf{0} \quad \text{in } \Body \\
        \StressTensor \BodyLateralNormal &= \mathbf{0} \quad \text{on } \BodyLateralSurface
    \end{align*}
    
    \item \textbf{Saint Venant Class:} The tractions on the end $\Section_L$ (at $x=L$) are statically equivalent to a transverse force $\Force$ applied at the centroid. We follow Ie\c{s}an's $(P_2)$ formulation, which sets the load at $x=0$. The resultant force on the section $\Section_0$ is $\Force$ and the resultant moment is $\mathbf{0}$.
    \begin{align}
        \int_{\Section_0} \StressTensor (-\FrameBasis) \, da &= \Force, \quad \text{with } \Force \cdot \FrameBasis = 0 \\
        \int_{\Section_0} \bm{r} \times (\StressTensor (-\FrameBasis)) \, da &= \mathbf{0} 
    \end{align}
    This implies the internal resultants at a section $x > 0$ are $\Force(x) = \Force$ and $\Moment(x) = x (\FrameBasis \times \Force)$.

    \item \textbf{Constitutive Law:}
    \begin{equation}
        \StressTensor = \lambda \operatorname{tr}(\StrainTensor) \One + 2\mu \StrainTensor
    \end{equation}

    \item \textbf{Ie\c{s}an's Postulate}
    The axial derivative of the displacement is rigid on each section:
    \begin{equation}
    \partial_x \bm{u}(x,\bm{r})=\bm{a}_1+\bm{\theta}_1\times\bm{r}\qquad(\bm{a}_1,\bm{\theta}_1\ \text{constant vectors}).
    \label{eq:iesan-postulate-2}
    \end{equation}
    Equivalently, stresses depend at most linearly on $\reflenx$ and the torsional amplitude vanishes.
\end{enumerate}

\subsubsection{Solution}

\begin{equation}
\begin{aligned}
\IntRoR\,\bm{c}_{\Section} + A\,\CentroidLocation\,c_{\varepsilon}&= \mathbf{0} \\
\bm{c}_{\Section}\cdot \CentroidLocation + c_{\varepsilon}& = 0 \\
\Jsv \,c_{\vartheta}&= -\FrameBasis\cdot\left(\int_{\Section} \bm{r}\times \nabla \bm{\WarpShearIesan}
+ \bm{W}_{(b)} \mathrm{~d}A \right)\bm{b}_{\Section} \\
&= -\int_{\Section} \Big[(\FrameBasis\times \bm{r})\cdot \nabla \WarpShearIesan
+ \tfrac{1}{2}\,\nu\, r^2\,(\FrameBasis\times \bm{r})\cdot \bm{b}_{\Section}\Big]\, da \\
\end{aligned} \\
\label{eq:iesan-2-22}
\end{equation}

Solving \eqref{eq:iesan-2-22} for \(c_{\varepsilon}\) furnishes:
\[
\begin{aligned}
c_{\vartheta}
&=- \frac{1}{\Jsv}\int_{\Section} \Big[(\FrameBasis\times \bm{r})\cdot \nabla \WarpShearIesan
+ \tfrac{1}{2}\,\nu\, r^2 \,(\FrameBasis\times \bm{r})\cdot \bm{b}_{\Section}\Big]\mathrm{~d}A \\
 &=- \frac{1}{\Jsv}\int_{\Section} \Big[(\FrameBasis\times \bm{r})\cdot (\nabla \bm{\WarpShearIesan}^T
+ \tfrac{1}{2}\,\nu\, r^2 \,\oneS) \bm{b}_{\Section}\Big]\mathrm{~d}A \\
 &= \FrameBasis\cdot\frac{1}{\Jsv}\left[\int_{\Section} \bm{r}\times (\nabla \bm{\WarpShearIesan}^T
+ \tfrac{1}{2}\,\nu\, r^2 \,\oneS)\mathrm{~d}A \right] \bm{b}_{\Section} \\
\end{aligned}
\]


From \eqref{eq:2-18} and \eqref{eq:2-22}:
\begin{equation*}
    \begin{aligned}
        c_{1,2,3} &= 0\\
        c_{\vartheta} &=- \frac{1}{\Jsv}\int_{\Section} \Big[(\FrameBasis\times \bm{r})\cdot \nabla (\bm{\WarpShearIesan}\cdot\bm{b}_{\Section})
+ \tfrac{1}{2}\,\nu\, r^2 \,(\FrameBasis\times \bm{r})\cdot \bm{b}_{\Section}\Big]\mathrm{~d}A \\
\bm{b}_{\Section} &= -\,E^{-1}\Big(\IntRoR - A\,\CentroidLocation\otimes\,\CentroidLocation\Big)^{-1}\ShearForce \\
 b_\varepsilon &=-\CentroidLocation\cdot\bm{b}_{\Section} \\
b_\vartheta &=0 \\
    \tilde{\bm{w}}_{\Section}&= \mathbf{0} \\
    \FrameBasis\cdot\tilde{\bm{w}} &= \WarpShearIesan \\
    \end{aligned}
\end{equation*}
where $r^2 = \bm{r} \cdot \bm{r}$, $\nu$ is Poisson's ratio, and $\WarpShearIesan$ is the out-of-plane warping function defined by \eqref{eq:bvp-warp-shear-iesan}.


\subsubsection{Warping}

\begin{equation}
\left\{\begin{aligned}
\Delta (\bm\WarpShearIesan\cdot\bm{b}_{\Section}) &= -\,2\big(\bm{r} -\CentroidLocation\big)\cdot\bm{b}_{\Section} \quad &\text{in } \Section \\
\partial_{\bm{\nu}} (\bm\WarpShearIesan\cdot\bm{b}_{\Section}) &= 
\bm{\nu}\cdot\nu\;\Big(\operatorname{dev}\bm{r}\otimes\bm{r}
 -\bm{r}\otimes\CentroidLocation\Big)\bm{b}_{\Section}
\quad &\text{on } \SectionBoundary \\
\int_{\Section}\bm\WarpShearIesan\cdot\bm{b}_{\Section}\,\mathrm{d}A &= 0
\end{aligned}\right.
\label{eq:bvp-warp-shear-iesan}
\end{equation}


\begin{equation*}
\left\{
\begin{aligned}
\Delta\varphi&=-2\big(\bm{b}\cdot\bm{r}+b_{\varepsilon}\big)\quad\text{in }\Omega,
\\
\partial_{\bm{\nu}}\varphi&=-\,\nu\,\bm{\pi}(\bm{r};\bm{b})\cdot\bm{\nu}\quad\text{on }\partial\Omega,
\\
\int_\Omega \varphi\,dA&=0.
\end{aligned}
% \label{eq:phiBVP}
\right.
\end{equation*}
%
\begin{Details}
or, more explicitly
%
\begin{equation*}
\left\{\begin{aligned}
\Delta \WarpShearIesan&= -\,2\big(\bm{r} -\CentroidLocation\big)\cdot\bm{b}_{\Section} \quad \text{in } \Section \\
\partial_{\bm{\nu}} \WarpShearIesan &= 
\nu\;\bm{b}_{\Section}\cdot\Big(\bm{r}\otimes\bm{r}
-\tfrac{1}{2}\, r^2\,\mathbf{1}
 -\CentroidLocation\otimes\bm{r}\Big)\cdot\bm{\nu}
\quad \text{on } \SectionBoundary .
\end{aligned}\right.
\end{equation*}
\end{Details}
%
For the BVP
\eqref{eq:bvp-warp-shear-iesan}, integrability of \emph{both} component
problems fixes \(\CentroidLocation\) at the \emph{area centroid}
\(\SectionAverage{\bm{r}}\).
%

\clearpage
\subsubsection{Displacement}
The displacement $\disp = u_x \FrameBasis + \disp_{\Section}$ (ignoring rigid body motion) takes the form:

\[
\begin{aligned}
\hat{\bm{u}}_{\Section}(x,\bm{r})
&= -\tfrac16\,x^{3}\,\bm{b}_{\Section}
+c_{\vartheta}\,x\,(\FrameBasis\times\bm{r})
+x\,\nu\Big(\tfrac12 r^{2}\,\bm{b}_{\Section}-(\bm{b}_{\Section}\cdot\bm{r})\,\bm{r}\Big)
-x\,\nu\,b_\varepsilon\,\bm{r}\,\\
\FrameBasis\cdot\hat{\bm{u}}(x\,\bm{r})
&=\tfrac12\,x^{2}\,(\bm{r}\cdot\bm{b}_{\Section} + b_{\varepsilon})
+c_{\vartheta}\,\WarpTwistIesan(\bm{r})+\WarpShearIesan(\bm{r})
\end{aligned}
\]
After eliminating \(b_{\varepsilon} = -\CentroidLocation\cdot\bm{b}_{\Section}\):
\begin{equation*}
\begin{aligned}
\hat{\bm{u}}_{\Section} \{\bm{b},\bm{c}\} &= -\tfrac{1}{6}\,x^3\,\bm{b}_{\Section}
%-\tfrac{1}{2}\,x^{\,2}\,\bm{c}_{\Section}
+ c_{\vartheta}\, x\,(\FrameBasis\times \bm{r})
 % +x \sum_{\beta} b_\beta\,\bm w^{(\beta)}(\bm{r})
 +x\nu\bigl(\tfrac{1}{2}r^2\mathbf{1}_{\Section} - \bm{r}\otimes\bm{r}\bigr)\bm{b}_{\Section}
 +x\nu\; \CentroidLocation\cdot\bm{b}_{\Section}\,\bm{r}
% + \tilde{\bm w}(\bm{r}) 
\\
\FrameBasis\cdot \hat{\bm{u}}(\reflenx,\bm{r})&=
~~\tfrac{1}{2}x^2\,\big(\bm{r}-\CentroidLocation\big)\cdot\bm{b}_{\Section}
%+ \bm{c}_{\Section}\cdot \bm{r}\,\reflenx
+ c_{\vartheta}\,\varphi(\bm{r}) + \WarpShearIesan(\bm{r}) .
\end{aligned}
\end{equation*}


Written in terms of the Poisson deviator  $\operatorname{dev}\bm{r}\otimes\bm{r} :=\bm{r}\otimes\bm{r}-\tfrac12 r^{2}\mathbf{1}$,

\begin{equation}
\boxed{
\begin{aligned}
\hat{\bm{u}}_{\Section} \{\bm{b},\bm{c}\} &=
 c_{\vartheta}\, x\,(\FrameBasis\times \bm{r})
-\tfrac{1}{6}\,x^3\,\bm{b}_{\Section}
 -x\nu \, \bigl(\operatorname{dev}\bm{r}\otimes\bm{r} - \bm{r}\otimes\CentroidLocation\bigr)\;\bm{b}_{\Section}
\\
%
\FrameBasis\cdot \hat{\bm{u}}\{\bm{b},\bm{c}\}&=
  c_{\vartheta}\,\WarpTwistIesan(\bm{r})
+ \WarpShearIesan(\bm{r})
+ \tfrac{1}{2}x^2\,\big(\bm{r} -\CentroidLocation\big)\cdot\bm{b}_{\Section}
\end{aligned}
}
\label{eq:iesan-displacement-2}
\end{equation}


\subsubsection{Strain}
\begin{equation*}
\StrainTensor(\hat{\bm{u}})= (\varepsilon \,\FrameBasis + \bm{\varepsilon}_{\gamma})\stackrel{\rm s}{\otimes}\FrameBasis - \nu\varepsilon \, \mathbf{1}_{\Section}
\quad\text{with}\quad
\begin{aligned}
\varepsilon &= (\FrameBasis\cdot\hat{\bm{u}})' \\
\bm\varepsilon_\gamma &= \hat{\bm{u}}_{\Section}'+\nabla_{\Section}(\FrameBasis \cdot\hat{\bm{u}})\\
\end{aligned}
\end{equation*}
where the components of \(\bm{\varepsilon}_{\gamma}\) are the engineering shear strains \(2\FrameBasis\cdot\bm{E}\,\FrameBasis_{\alpha}\)

\begin{Quote}
\noindent\emph{(iv) The $\bm{b}_{\Section}$-block.} Collect
\[
\begin{aligned}
\hat{\bm{u}}^{(b)}&=
-\tfrac16 x^3\,\bm{b}_{\Section}
\;-\;x\nu\big(\operatorname{dev}\bm{r}\otimes\bm{r} - \bm{r}\otimes\CentroidLocation\big)\bm{b}_{\Section}\\
&\qquad
+\tfrac12 x^2\,\FrameBasis\otimes(\bm{r}-\CentroidLocation)\,\bm{b}_{\Section}
\;+\;\FrameBasis\otimes\bm{\WarpShearIesan}(\bm{r})\,\bm{b}_{\Section}.
\end{aligned}
\]
Using
\[
\nabla \big[(\operatorname{dev}\bm{r}\otimes\bm{r}-\bm{r}\otimes\CentroidLocation) \, \bm{b}_{\Section}\big]
=\mathbf{1}_\Section\,((\bm{r} - \CentroidLocation)\cdot\bm{b}_{\Section}),
\]
and the cancellation
\(-\tfrac12x^2\FrameBasis\otimes\bm{b}_{\Section}+\tfrac12x^2\FrameBasis\otimes\bm{b}_{\Section}=\mathbf{0}\),
we obtain
\[
\begin{aligned}
\nabla\hat{\bm{u}}^{(b)}
&=\ \boxed{\,x\,(\FrameBasis\otimes\FrameBasis)\,((\bm{r}-\CentroidLocation)\cdot\bm{b}_{\Section})\,}
\;-\;x\nu\,((\bm{r}-\CentroidLocation)\cdot\bm{b}_{\Section})\,\mathbf{1}_\Section
\;+\;\FrameBasis\otimes \nabla(\bm{\WarpShearIesan}\cdot\bm{b}_{\Section})
\;-\;\nu\,\FrameBasis\otimes\bm{\pi}^{(b)}
\end{aligned}
\]
where \(\bm{\pi}^{(b)}:=(\operatorname{dev}(\bm{r}\otimes\bm{r})-\bm{r}\otimes\CentroidLocation)\bm{b}_{\Section}\).
\end{Quote}

\begin{Quote}
\medskip
\noindent\emph{(iv) The \(\bm{b}_{\Section}\)-block.}  Collect
\[
\begin{aligned}
\hat{\bm{u}}^{(b)}&=
-\tfrac16 x^3\,\bm{b}_{\Section}
\;-\;x\nu\big(\operatorname{dev}\bm{r}\otimes\bm{r} - \bm{r}\otimes\CentroidLocation\big)\bm{b}_{\Section}\\
&\qquad
+\tfrac12 x^2\,\FrameBasis\otimes(\bm{r}-\CentroidLocation)\,\bm{b}_{\Section}
\;+\;\FrameBasis\otimes\bm{\WarpShearIesan}(\bm{r})\,\bm{b}_{\Section}.
\end{aligned}
\]
Differentiating and using
\[
\nabla \big[(\operatorname{dev}\bm{r}\otimes\bm{r}-\bm{r}\otimes\CentroidLocation)\bm{b}_{\Section}\big]
=\mathbf{1}_\Section (\bm{r}-\CentroidLocation)\cdot\bm{b}_{\Section},
\]
together with the cancellation
\(-\tfrac12x^2\FrameBasis\otimes\bm{b}_{\Section}+\tfrac12x^2\FrameBasis\otimes\bm{b}_{\Section}=\mathbf{0}\),
gives
\[
\begin{aligned}
\nabla\hat{\bm{u}}^{(b)}
&=\ (\FrameBasis\otimes\FrameBasis)\,((\bm{r}-\CentroidLocation)\cdot\bm{b}_{\Section})
\;-\;x\nu\,((\bm{r}-\CentroidLocation)\cdot\bm{b}_{\Section})\,\mathbf{1}_\Section
% \\
% &\qquad
+ \FrameBasis\otimes\Big(\nabla \bm{\WarpShearIesan}\cdot\bm{b}_{\Section}\Big)
\;-\;\nu\,\FrameBasis\otimes\bm{\pi}^{(b)}
\end{aligned}
\]
where \(\bm{\pi}^{(b)}:=(\operatorname{dev}(\bm{r}\otimes\bm{r})-\bm{r}\otimes\CentroidLocation)\bm{b}_{\Section}\).
Hence
\[
\bm{E}^{(b)}
=((\bm{r}-\CentroidLocation)\cdot\bm{b}_{\Section})\,(\FrameBasis\otimes\FrameBasis)
\;-\;x\nu\,((\bm{r}-\CentroidLocation)\cdot\bm{b}_{\Section})\,\mathbf{1}_\Section
\;+\;\tfrac12\Big[\FrameBasis\otimes\bm{\varepsilon}^{(b)}+\bm{\varepsilon}^{(b)}\otimes\FrameBasis\Big],
\]
with \(\displaystyle \bm{\varepsilon}^{(b)}:=\nabla (\bm{\WarpShearIesan}\cdot\bm{b}_{\Section})-\nu\,\bm{\pi}^{(b)}\).
\end{Quote}

\begin{Details}
\begin{equation*}
\begin{aligned}
\bm E(\hat{\bm{u}})&= (\varepsilon \,\FrameBasis + \bm{\varepsilon}_{\gamma})\stackrel{\rm s}{\otimes}\FrameBasis \\
\varepsilon &= (\FrameBasis\cdot\hat{\bm{u}})' \\
&= x\,(\bm{r}\cdot\bm{b}_{\Section} + b_{\varepsilon}) \\ % ok
&= x\,(\bm{r}-\CentroidLocation)\cdot\bm{b}_{\Section},\\[2pt]
\bm\varepsilon_\gamma &= \hat{\bm{u}}_{\Section}'+\nabla_{\Section}(\FrameBasis \cdot\hat{\bm{u}})\\
&=\ c_{\vartheta}\big(\nabla_{\Section}\varphi-\FrameBasis\times\bm{r}\big)
+\nu\Big(\tfrac12 r^{2}\,\bm{b}_{\Section}-(\bm{b}_{\Section}\cdot\bm{r})\,\bm{r}\Big)
-\nu\,b_\varepsilon\,\bm{r}
+\nabla \WarpShearIesan \\
&=\ c_{\vartheta}\big(\nabla_{\Section}\varphi-\FrameBasis\times\bm{r}\big)
-\nu\operatorname{dev}(\bm{r}\otimes\bm{r})\,\bm{b}_{\Section}
-\nu\,b_\varepsilon\,\bm{r}
+\nabla \WarpShearIesan \\
&=\ c_{\vartheta}\big(\nabla_{\Section}\varphi-\FrameBasis\times\bm{r}\big)
-\nu\operatorname{dev}(\bm{r}\otimes\bm{r})\,\bm{b}_{\Section}
+\nu\,\bm{r}\otimes\CentroidLocation\;\bm{b}_{\Section}
+\nabla \WarpShearIesan \\
&=\ c_{\vartheta}\big(\nabla_{\Section}\varphi-\FrameBasis\times\bm{r}\big)
-\nu\bigl(\operatorname{dev}(\bm{r}\otimes\bm{r})-\bm{r}\otimes\CentroidLocation \bigr)\,\bm{b}_{\Section}
+\nabla (\bm{\WarpShearIesan}\cdot\bm{b}_{\Section}) \\
\end{aligned}
\end{equation*}

\begin{equation*}
\begin{aligned}
\bm E(\hat{\bm{u}})&= (\varepsilon \,\FrameBasis + \bm{\varepsilon}_{\gamma})\stackrel{\rm s}{\otimes}\FrameBasis \\
\varepsilon &= (\FrameBasis\cdot\hat{\bm{u}})' \\
&= x\,(\bm{r}\cdot\bm{b}_{\Section} + b_{\varepsilon}) \\ % ok
&= x\,(\bm{r}-\CentroidLocation)\cdot\bm{b}_{\Section},\\[2pt]
\bm\varepsilon_\gamma &= \hat{\bm{u}}_{\Section}'+\nabla_{\Section}(\FrameBasis \cdot\hat{\bm{u}})\\
&=\ c_{\vartheta}\big(\nabla \WarpTwistIesan + \FrameBasis\times\bm{r}\big)
+\nu\Big(\tfrac12 r^{2}\,\bm{b}_{\Section}-(\bm{b}_{\Section}\cdot\bm{r})\,\bm{r}\Big)
-\nu\,b_\varepsilon\,\bm{r}
+\nabla \WarpShearIesan \\
&=\ c_{\vartheta}\big(\nabla \WarpTwistIesan + \FrameBasis\times\bm{r}\big)
-\nu\operatorname{dev}(\bm{r}\otimes\bm{r})\,\bm{b}_{\Section}
-\nu\,b_\varepsilon\,\bm{r}
+\nabla \WarpShearIesan \\
&=\ c_{\vartheta}\big(\nabla \WarpTwistIesan + \FrameBasis\times\bm{r}\big)
-\nu\operatorname{dev}(\bm{r}\otimes\bm{r})\,\bm{b}_{\Section}
+\nu\,\bm{r}\otimes\CentroidLocation\;\bm{b}_{\Section}
+\nabla \WarpShearIesan \\
&=\ c_{\vartheta}\big(\nabla \WarpTwistIesan + \FrameBasis\times\bm{r}\big)
-\nu\bigl(\operatorname{dev}(\bm{r}\otimes\bm{r})-\bm{r}\otimes\CentroidLocation \bigr)\,\bm{b}_{\Section}
+\nabla (\bm{\WarpShearIesan}\cdot\bm{b}_{\Section}) \\
\end{aligned}
\end{equation*}
\end{Details}

\begin{equation}
\boxed{
\begin{aligned}
\bm E(\hat{\bm{u}})&= (\varepsilon \,\FrameBasis + \bm{\varepsilon}_{\gamma})\stackrel{\rm s}{\otimes}\FrameBasis \\
%
\varepsilon %&= (\FrameBasis\cdot\hat{\bm{u}})' \\
% &= x\,(\bm{r}\cdot\bm{b}_{\Section} + b_{\varepsilon}) \\ % ok
&= x\,(\bm{r}-\CentroidLocation)\cdot\bm{b}_{\Section},\\[2pt]
\bm\varepsilon_\gamma 
&=\ c_{\vartheta}\big(\nabla \WarpTwistIesan + \FrameBasis\times\bm{r}\big)
-\nu\bigl(\operatorname{dev}(\bm{r}\otimes\bm{r})-\bm{r}\otimes\CentroidLocation \bigr)\,\bm{b}_{\Section}
+\nabla (\bm{\WarpShearIesan}\cdot\bm{b}_{\Section}) \\
\end{aligned}}
\label{eq:iesan-strain-2}
\end{equation}
Let
\[
\bm{\varepsilon}^{(b)}:=\nabla (\bm{\WarpShearIesan}\cdot\bm{b}_{\Section})
-\nu\,\big(\operatorname{dev}(\bm{r}\otimes\bm{r})-\bm{r}\otimes\CentroidLocation\big)\bm{b}_{\Section},
\]
so that
\begin{equation}
\StrainTensor^{(b)}
=((\bm{r}-\CentroidLocation)\cdot\bm{b}_{\Section})\,\FrameBasis\otimes\FrameBasis
-x\nu\,((\bm{r}-\CentroidLocation)\cdot\bm{b}_{\Section})\,\mathbf{1}_\Section
+\tfrac12\Big(\FrameBasis\otimes\bm{\varepsilon}^{(b)}+\bm{\varepsilon}^{(b)}\otimes\FrameBasis\Big),
\label{eq:iesan-strain-tensor-b}
\end{equation}

\subsubsection{Stress}
Using \eqref{eq:iesan-strain-2} with the constitutive equation furnishes:

\begin{equation}
\boxed{%
\begin{aligned}
\ShearStress(\hat{\bm{u}})&=
G\Big[c_{\vartheta}\big(\nabla \WarpTwistIesan + \FrameBasis\times \bm{r}\big)
- \nu\,\bm{r}\,\big(\bm{r} -\CentroidLocation\big)\cdot\bm{b}_{\Section}
+ \tfrac{1}{2}\,\nu\,r^2\,\bm{b}_{\Section}
+ \nabla_{\Section} \psi\Big] \\
&=G\Big[\ c_{\vartheta}\big(\nabla\WarpTwistIesan+\FrameBasis\times\bm{r}\big)
-\nu\bigl(\operatorname{dev}(\bm{r}\otimes\bm{r})-\bm{r}\otimes\CentroidLocation \bigr)\,\bm{b}_{\Section}
+\nabla (\bm{\WarpShearIesan}\cdot\bm{b}_{\Section})\Big] \\
\AxialStress(\hat{\bm{u}}) &=
\reflenx E \big(\bm{r} -\CentroidLocation\big)\cdot\bm{b}_{\Section}.
\end{aligned}}
\label{eq:iesan-stress-2}
\end{equation}
\begin{Quote}
\medskip
\noindent\emph{(iv) The \(\bm{b}_{\Section}\)-block.}

From \eqref{eq:iesan-strain-tensor-b}, with \(\operatorname{tr}\StrainTensor^{(b)}=(1-2\nu x)\,((\bm{r}-\CentroidLocation)\cdot\bm{b}_{\Section})\).
Therefore
\begin{equation}
\begin{aligned}
\StressTensor^{(b)}
&=2G\,\StrainTensor^{(b)}+\lambda\,\operatorname{tr}(\StrainTensor^{(b)})\,\mathbf{1} \\
&=\ \Big[\,2G+\lambda(1-2\nu x)\,\Big]\;((\bm{r}-\CentroidLocation)\cdot\bm{b}_{\Section})\ \FrameBasis\otimes\FrameBasis\\
&\quad+\ \Big[\,\lambda(1-2\nu x)-2G\,\nu x\,\Big]\;((\bm{r}-\CentroidLocation)\cdot\bm{b}_{\Section})\ \mathbf{1}_\Section\\
&\quad+\ G\Big[\FrameBasis\otimes\bm{\varepsilon}^{(b)}+\bm{\varepsilon}^{(b)}\otimes\FrameBasis\Big].
\end{aligned}
\label{eq:iesan-stress-tensor-flexure-2}
\end{equation}
\end{Quote}

%


\subsubsection{Observations}
\begin{itemize}
    \item The solution presented is exact for the relaxed problem (Condition 4) under the Saint-Venant hypothesis (Condition 5). By Saint-Venant's principle, it serves as a globally accurate approximation for a long cylinder (where $L$ is large compared to the dimensions of $\Section$) under any system of end tractions that is statically equivalent to $\Force$.
    
    \item The linearity of the axial stress $\AxialStress$ in $\bm{r}$ (i.e., $\sigma(\bm{r})' \propto {\bm{b}_{\Section}} \cdot \bm{r}$) is a crucial and necessary condition for the Saint-Venant hypothesis ($T_{yy}=T_{zz}=T_{yz}=0$) to be compatible with the equations of elasticity.
\end{itemize}


\subsection{General Solution}
\subsubsection{Displacement}

The complete Saint Venant solution for flexure-bending-torsion-extension is parameterized by \(\bm{a} = (\IesanLoadM, a_{\varepsilon}, a_{\vartheta})\in\mathbb{R}^4\) and \(\bm{b}_{\Section}\in \mathbb{R}^2\).
Combining \eqref{eq:iesan-displacement-1} and \eqref{eq:iesan-displacement-2} furnishes:

\begin{equation}
\begin{aligned}
  \hat{\bm{u}}\{\bm{a},\bm{b}_{\Section}\}
    &= 
      \left(-\tfrac{1}{2}x^2 \mathbf{1} - \nu\operatorname{dev}\bm{r}\otimes\bm{r} + x\, \FrameBasis\otimes\bm{r}\right)\IesanLoadM \\ 
    &+ (x\FrameBasis - \nu \bm{r})a_{\varepsilon} \\
    &+(x\; \FrameBasis\times \bm{r} + \varphi \FrameBasis)a_{\vartheta}\\
    & + (x\,\FrameBasis\times\bm{r}+ \WarpTwistIesan\FrameBasis)\, c_{\vartheta} \\
    &-\tfrac{1}{6}\,x^3\,\bm{b}_{\Section}
      - x\nu \, \bigl(\operatorname{dev}\bm{r}\otimes\bm{r} - \bm{r}\otimes\CentroidLocation\bigr)\;\bm{b}_{\Section} \\
    &+\tfrac{1}{2}x^2\,\FrameBasis\otimes\big(\bm{r} -\CentroidLocation\big)\bm{b}_{\Section}
    + \FrameBasis\otimes\bm{\WarpShearIesan}(\bm{r})\;\bm{b}_{\Section} .
\end{aligned}
% \label{eq:iesan-general-solution}
\end{equation}

\begin{equation}
\begin{aligned}
  \hat{\bm{u}}\{\bm{a},\bm{b}_{\Section}\}
    &= 
      \left(-\tfrac{1}{2}x^2 \mathbf{1} +\bm{W}_{(a)} + x\, \FrameBasis\otimes\bm{r}\right)\IesanLoadM \\ 
    &+ (x\FrameBasis - \nu \bm{r})a_{\varepsilon} \\
    &+(x\; \FrameBasis\times \bm{r} + \varphi \FrameBasis) (a_{\vartheta} + c_{\vartheta})\\
    &-\tfrac{1}{6}\,x^3\,\bm{b}_{\Section}
      + x \bm{W}_{(b)}\;\bm{b}_{\Section} \\
    &+\tfrac{1}{2}x^2\,\FrameBasis\otimes\big(\bm{r} -\CentroidLocation\big)\bm{b}_{\Section}
    + \FrameBasis\otimes\bm{\WarpShearIesan}(\bm{r})\;\bm{b}_{\Section} .
\end{aligned}
% \label{eq:iesan-general-solution}
\end{equation}
\begin{Details}
\[
\begin{aligned} \hat{\bm{u}}\left\{\bm{a}, \bm{b}_{\Section}\right\}= 
& \underbrace{\left(-\tfrac{1}{2} x^2 \mathbf{1}\right) a_{\Section}}_{\text {axis (transverse) }}
-\underbrace{\nu (\operatorname{dev} \bm{r} \otimes \bm{r})\, \IesanLoadM}_{\text {plane Poisson}}
+\underbrace{(x \FrameBasis \otimes r) \IesanLoadM}_{\text {axial linear in } r} \\ 
& +\underbrace{(x \FrameBasis)\, a_{\varepsilon}}_{\text {axial extension}}
-\underbrace{\nu \,\bm{r}\, a_{\varepsilon}}_{\text {in-plane Poisson }} \\ 
& +\underbrace{(x \FrameBasis \times \bm{r}) a_{\vartheta}}_{\text {rigid twist }}
+\underbrace{\varphi(r) \FrameBasis (a_{\vartheta} + c_{\vartheta})}_{\text {torsion warping }} \\ 
& -\underbrace{\left.\tfrac{1}{6} x^3 \,\bm{b}_{\Section}\right.}_{\text {axis (transverse) }}
+\underbrace{\left(c_{\vartheta} x\, \FrameBasis \times \bm{r}\right)}_{\text {rigid twist from } b_{\Section}} \\ 
& -\underbrace{\left. x \nu\left(\text { dev } \bm{r}\otimes\bm{r} - \bm{r} \otimes \CentroidLocation\right)\right. \bm{b}_{\Section}}_{\text {in-plane Poisson }} \\ 
& +\underbrace{\frac{1}{2} x^2 \FrameBasis \otimes\left(\bm{r}-\CentroidLocation\right) \bm{b}_{\Section}}_{\text {axial linear in } \bm{r}}
+\underbrace{\FrameBasis \otimes \bm{\WarpShearIesan}(\bm{r}) \bm{b}_{\Section}}_{\text {flexural warping }}
\end{aligned}
\]
\end{Details}

Grouping by parameters, the kinematics and amplitudes reproduce each Saint–Venant block:
\[
\begin{aligned}
\IesanLoadM:&\quad
-\tfrac{1}{2} x^2\,\IesanLoadM\;-\;\nu\,\operatorname{dev}\bm{r}\otimes\bm{r}\,\IesanLoadM
\;+\;x\,\FrameBasis\otimes\bm{r}\,\IesanLoadM,\\
a_\varepsilon:&\quad x\,\FrameBasis\,a_\varepsilon\;-\;\nu\,\bm{r}\,a_\varepsilon,\\
a_\vartheta:&\quad x\,\FrameBasis\times\bm{r}\,a_\vartheta\;+\;\varphi(\bm{r})\,\FrameBasis\,a_\vartheta,\\
\bm{b}_{\Section}:&\quad
-\tfrac16 x^3\,\bm{b}_{\Section}\;-\;x\,\nu\big(\operatorname{dev}\bm{r}\otimes\bm{r}-\bm{r}\otimes\CentroidLocation\big)\bm{b}_{\Section}
\;+\;\tfrac{1}{2} x^2\,\FrameBasis\otimes(\bm{r}-\CentroidLocation)\,\bm{b}_{\Section}
\;+\;\FrameBasis\otimes\bm\psi(\bm{r})\,\bm{b}_{\Section},\\
c_\vartheta:&\quad \,c_\vartheta\,x\,\FrameBasis\times\bm{r}\;+\;c_\vartheta\,\varphi(\bm{r})\,\FrameBasis.
\end{aligned}
\]
%
This decomposes into longitudinal \(u_x\) and transverse \(\hat{\bm{u}}_{\Section}\) components:
\begin{equation*}
\begin{aligned}
\hat{\bm{u}}\{\bm{a},\bm{b}_{\Section}\} &= \;...
\end{aligned}
\end{equation*}

\subsubsection{Strain}


\paragraph{Gradient.}
Decompose $\nabla \hat{\bm{u}} = \FrameBasis\otimes\partial_x\hat{\bm{u}} + \gradS\hat{\bm{u}}$.


\emph{Axial part} $\FrameBasis\otimes\partial_x\hat{\bm{u}}$:
\[
\begin{aligned}
\partial_x\hat{\bm{u}}
&=\; \Big(-x\,\one + \FrameBasis\otimes\bm{r}\Big)\IesanLoadM
\;+\; \FrameBasis\,a_\varepsilon
\;+\; (\FrameBasis\times\bm{r})\,(a_\vartheta+c_\vartheta)
\\&\quad
-\tfrac12 x^2\,\bm{b}_{\Section}
\;-\;\nu\bigl(\dev(\bm{r}\otimes\bm{r}) - \bm{r}\otimes\CentroidLocation\bigr)\bm{b}_{\Section}
\;+\; x\,\FrameBasis\otimes(\bm{r}-\CentroidLocation)\bm{b}_{\Section} .
\end{aligned}
\]

\emph{Transverse part} $\gradS\hat{\bm{u}}$:
\[
\begin{aligned}
\gradS\hat{\bm{u}}
&=\; \Big[-\nu\,(\bm{r}\cdot\IesanLoadM)\,\oneS \;+\; x\,\FrameBasis\otimes \IesanLoadM\Big]
\;-\;\nu\,a_\varepsilon\,\oneS
\\&\quad
+\; x\,( \FrameBasis\times \oneS)\,a_\vartheta
\;+\; \FrameBasis\otimes(\gradS\varphi)\,a_\vartheta
\;+\; x\,( \FrameBasis\times \oneS)\,c_\vartheta
\\&\quad
-\; x\nu\,\bigl[(\bm{r}-\CentroidLocation)\cdot\bm{b}_{\Section}\bigr]\oneS
\;+\; \tfrac12 x^2\,\FrameBasis\otimes \bm{b}_{\Section}
\;+\; \FrameBasis\otimes\bigl(c_\vartheta\,\gradS\WarpTwistIesan\bigr)
\\&\quad
+\; \FrameBasis\otimes \gradS\!\big(\bm{\WarpShearIesan}\cdot \bm{b}_{\Section}\big).
\end{aligned}
\]

Combine these to write
\[
\begin{aligned}
\nabla \hat{\bm{u}}
&= (\FrameBasis\otimes\FrameBasis)\,\Big[a_\varepsilon+(\bm{r}\cdot\IesanLoadM)
+ x\,\big((\bm{r}-\CentroidLocation)\cdot\bm{b}_{\Section}\big)\Big]
\\&\quad
+\FrameBasis\otimes\Big\{
-x\,\IesanLoadM
+ (\FrameBasis\times\bm{r})\,(a_\vartheta+c_\vartheta)
-\tfrac12 x^2\,\bm{b}_{\Section}
% \\&\qquad\qquad
-\nu\bigl(\dev(\bm{r}\otimes\bm{r})-\bm{r}\otimes\CentroidLocation\bigr)\bm{b}_{\Section}
\Big\}
\\&\quad
+ \FrameBasis\otimes\Big\{
~x\,\IesanLoadM + \tfrac12 x^2\,\bm{b}_{\Section}
+ (\gradS\varphi)\,a_\vartheta
+ c_\vartheta\,\gradS\WarpTwistIesan
+ \gradS \big(\bm{\WarpShearIesan}\cdot\bm{b}_{\Section}\big)
\Big\}
\\&\quad
-\nu\Big[a_\varepsilon+(\bm{r}\cdot\IesanLoadM)
+ x\,\big((\bm{r}-\CentroidLocation)\cdot\bm{b}_{\Section}\big)\Big]\oneS
\\&\quad
+ x\,(\FrameBasis\times \oneS)\,(a_\vartheta+c_\vartheta),
\end{aligned}
\]
where the last line is planar-skew (pure rotation).

%----------------------------------------------------------------------
% Infinitesimal strain (symmetric gradient)
%----------------------------------------------------------------------

\paragraph{Infinitesimal strain.}
\[
\StrainTensor \;\equiv\; \sym\nabla\hat{\bm{u}}
\;=\; \tfrac12\big(\nabla\hat{\bm{u}}+(\nabla\hat{\bm{u}})^{\mathrm t}\big).
\]
Using that $\sym(\FrameBasis\times\oneS)=\mathbf{0}$  and collecting symmetric parts:

\[
\begin{aligned}
\StrainTensor
&= \underbrace{\Big[a_\varepsilon+(\bm{r}\cdot\IesanLoadM)
+ x\,\big((\bm{r}-\CentroidLocation)\cdot\bm{b}_{\Section}\big)\Big]}_{\displaystyle \varepsilon}
\,(\FrameBasis\otimes\FrameBasis)
\\[2pt]
&\quad
+\;\sym\Big\{\FrameBasis\otimes \bm V_\Sigma\Big\}
\;-\;
\nu\,\underbrace{\Big[a_\varepsilon+(\bm{r}\cdot\IesanLoadM)
+ x\,\big((\bm{r}-\CentroidLocation)\cdot\bm{b}_{\Section}\big)\Big]}_{\displaystyle \varepsilon}
\,\oneS ,
\end{aligned}
\]
with the mixed (axial–transverse) vector
\[
\bm V_\Sigma
= (\FrameBasis\times\bm{r})\,(a_\vartheta+c_\vartheta)
\;-\;\nu\bigl(\dev(\bm{r}\otimes\bm{r})-\bm{r}\otimes\CentroidLocation\bigr)\bm{b}_{\Section}
\;+\; (\nabla\WarpTwistIesan)\,a_\vartheta
\;+\; c_\vartheta\,\gradS\WarpTwistIesan
\;+\; \nabla\!\big(\bm{\WarpShearIesan}\cdot\bm{b}_{\Section}\big) .
\]
The infinitesimal strain is
\[
\StrainTensor
= \varepsilon\,\FrameBasis\otimes\FrameBasis
\;+\;\sym\{\FrameBasis\otimes \bm V_\Sigma\}
\;-\;\nu\,\varepsilon\,\oneS,
\]
where
\begin{subequations}
\begin{align}
\PointAxialStrain
&= a_\varepsilon + \bm{r}\cdot\IesanLoadM
+ x\,(\bm{r}-\CentroidLocation)\cdot\bm{b}_{\Section},
 \label{eq:iesan-point-axial-strain} \\
\bm V_\Sigma
&= (\FrameBasis\times\bm{r} + \gradS \WarpTwistIesan)\,(a_\vartheta+c_\vartheta)
\;+\; \gradS \big(\bm{\WarpShearIesan}\cdot\bm{b}_{\Section}\big)
\;-\;\nu\bigl(\dev(\bm{r}\otimes\bm{r})-\bm{r}\otimes\CentroidLocation\bigr)\bm{b}_{\Section}\\
\PointShearStrain &= (a_{\vartheta} + c_{\vartheta})(\FrameBasis\times \bm{r} + \nabla\WarpTwistIesan) - \nu \bigl(\operatorname{dev}\bm{r}\otimes\bm{r}
- \bm{r}\otimes\CentroidLocation\bigr)\bm{b}_{\Section}
+ (\nabla \bm{\WarpShearIesan})^T \bm{b}_{\Section} \label{eq:iesan-point-shear-strain}
\end{align}
\label{eq:iesan-point-strain}
\end{subequations}



\begin{Quote}
\textbf{Complete infinitesimal strain.}

Summing the four blocks,
\[
\boxed{
\begin{aligned}
\bm{E}
&=\ (\bm{r}\cdot\IesanLoadM\;+\;a_\varepsilon)\,\big(\FrameBasis\otimes\FrameBasis-\nu\,\mathbf{1}_\Section\big) \\[2pt]
&\quad
+\ \tfrac12\Big(\FrameBasis\otimes\bm{\varepsilon}_{\vartheta}+\bm{\varepsilon}_{\vartheta}\otimes\FrameBasis\Big)(a_\vartheta+c_\vartheta)
+\tfrac12\Big(\FrameBasis\otimes\bm{\varepsilon}^{(b)}+\bm{\varepsilon}^{(b)}\otimes\FrameBasis\Big)\\[2pt]
&\quad
+\ ((\bm{r}-\CentroidLocation)\cdot\bm{b}_{\Section})\,(\FrameBasis\otimes\FrameBasis)
\;-\;x\nu\,((\bm{r}-\CentroidLocation)\cdot\bm{b}_{\Section})\,\mathbf{1}_\Section
\end{aligned}}
\]
with
\[
\bm{\varepsilon}_{\vartheta}= \,\FrameBasis\times\bm{r}
+\nabla \WarpTwistIesan
\qquad
\bm{\varepsilon}^{(b)}=\nabla (\bm{\WarpShearIesan}\cdot\bm{b}_{\Section})
-\nu\,\big(\operatorname{dev}\bm{r}\otimes\bm{r}-\bm{r}\otimes\CentroidLocation\big)\bm{b}_{\Section}.
\]
\end{Quote}

\begin{Quote}

\[
\boxed{
\begin{aligned}
\bm{E}
&=\ (\bm{r}\cdot\IesanLoadM + a_{\varepsilon})\,\big(\FrameBasis\otimes\FrameBasis-\nu\,\mathbf{1}_\Section\big)\\[2pt]
&\quad
+\ \tfrac12\Big[\FrameBasis\otimes\bm{\varepsilon}_{\vartheta}+\bm{\varepsilon}_{\vartheta}\otimes\FrameBasis\Big](a_\vartheta+c_\vartheta)
\;+\;\tfrac12 \Big(\FrameBasis\otimes\bm{\varepsilon}^{(b)}+\bm{\varepsilon}^{(b)}\otimes\FrameBasis\Big)\\[2pt]
&\quad
+\ x\,((\bm{r}-\CentroidLocation)\cdot\bm{b}_{\Section})\,(\FrameBasis\otimes\FrameBasis)
\;-\;x\nu\,((\bm{r}-\CentroidLocation)\cdot\bm{b}_{\Section})\,\mathbf{1}_\Section
\end{aligned}}
\]
\[
\bm{\varepsilon}_{\vartheta}= \,\FrameBasis\times\bm{r}
+\nabla \varphi,
\qquad
\bm{\varepsilon}^{(b)}=\nabla\big(\bm{\WarpShearIesan}\cdot\bm{b}_{\Section}\big)
-\nu\,\big(\operatorname{dev}(\bm{r}\otimes\bm{r})-\bm{r}\otimes\CentroidLocation\big)\bm{b}_{\Section}.
\]
\end{Quote}

\paragraph{Remarks}
\begin{itemize}
\item The axial normal strain is \(
\varepsilon=a_\varepsilon+ \bm{r}\cdot\IesanLoadM
+ x\,\big((\bm{r}-\CentroidLocation)\cdot\bm{b}_{\Section}\big)
\).
\item The planar normal strains are isotropic in the section and equal to
\(
-\nu\,\varepsilon\,\oneS
\).
\item The only contributors to symmetric axial–shear components are the mixed vector
$\bm V_\Sigma$; all $x(\FrameBasis\times\oneS)$ terms are planar-skew and thus drop out of
$\StrainTensor$.
\end{itemize}

\subsubsection{Stress}

For linear elastic isotropy in 3D,
\[
\StressTensor=2G\,\StrainTensor
  +\lambda\,\operatorname{tr}(\StrainTensor)\,\mathbf{1}.
\]

\begin{Quote}
% \mathbf{1}=\FrameBasis\otimes\FrameBasis+\mathbf{1}_\Section, 
% \mathbf{1}_\Section=\FrameBasis_y\otimes\FrameBasis_y+\FrameBasis_z\otimes\FrameBasis_z.
% Linear elastic isotropy: \bm{\sigma}=2G\,\bm{\varepsilon}+\lambda\,\operatorname{tr}(\bm{\varepsilon})\,\mathbf{1}
% with G the shear modulus and \lambda the Lamé parameter
% \big(\lambda=\tfrac{2G\nu}{1-2\nu}=\tfrac{E\nu}{(1+\nu)(1-2\nu)}\big).

From the strain derived previously,
\[
\operatorname{tr} \StrainTensor
=(1-2\nu)\big(\bm{r}\cdot\bm{a}_\Section+a_\varepsilon\big)
+\big(1-2\nu x\big)\,(\bm{r}-\CentroidLocation)\cdot\bm{b}_\Section,
\]
since \(\operatorname{tr}\FrameBasis\otimes\bm v=\FrameBasis\cdot\bm v=0\) for any in–section \(\bm v\).

\textbf{Stress decomposition by blocks.}
Write \(\bm{\sigma}=\bm{\sigma}^{(a)}+\bm{\sigma}^{(\varepsilon)}+\bm{\sigma}^{(\vartheta,c_\vartheta)}+\bm{\sigma}^{(b)}\) corresponding to the four parts of \(\StrainTensor\). 

\medskip
\noindent\emph{(iii) The \((a_\vartheta,c_\vartheta)\)-block.}
Here
\[
\StrainTensor^{(\vartheta,c_\vartheta)}
=\tfrac12\Big[\FrameBasis\otimes\bm{\varepsilon}_{\vartheta}+\bm{\varepsilon}_{\vartheta}\otimes\FrameBasis\Big],
\qquad
\bm{\varepsilon}_{\vartheta}:=(a_\vartheta-c_\vartheta)\,(\FrameBasis\times\bm{r})
+a_\vartheta\,\nabla \WarpTwistIesan
+c_\vartheta\,\nabla_\Section\WarpTwistIesan,
\]
with \(\bm{\varepsilon}_{\vartheta}\) in–section and \(\operatorname{tr}\bm{E}^{(\vartheta,c_\vartheta)}=0\). 
Hence
\[
\boxed{\quad
\StressTensor^{(\vartheta,c_\vartheta)}
=2G\,\bm{E}^{(\vartheta,c_\vartheta)}
=G\Big[\FrameBasis\otimes\bm{\varepsilon}_{\vartheta}+\bm{\varepsilon}_{\vartheta}\otimes\FrameBasis\Big]. \quad}
\]
\end{Quote}


\begin{Quote}
\textbf{Complete Cauchy stress.}
Summing the terms \eqref{eq:iesan-stress-tensor-flexure-2}, ...,
\[
\boxed{
\begin{aligned}
\StressTensor
&=\ 2G(1+\nu)\,\big(\bm{r}\cdot\IesanLoadM+a_\varepsilon\big)\ \FrameBasis\otimes\FrameBasis
\;+\;G\Big[\FrameBasis\otimes\bm{\varepsilon}_{\vartheta}+\bm{\varepsilon}_{\vartheta}\otimes\FrameBasis\Big]\\[2pt]
&\quad+\ \Big[\,2G+\lambda(1-2\nu x)\,\Big]\;((\bm{r}-\CentroidLocation)\cdot\bm{b}_{\Section})\ \FrameBasis\otimes\FrameBasis\\[2pt]
&\quad+\ \Big[\,\lambda(1-2\nu x)-2G\,\nu x\,\Big]\;((\bm{r}-\CentroidLocation)\cdot\bm{b}_{\Section})\ \mathbf{1}_\Section \\[2pt]
&\quad+\ G\Big[\FrameBasis\otimes\bm{\varepsilon}^{(b)}+\bm{\varepsilon}^{(b)}\otimes\FrameBasis\Big],
\end{aligned}}
\]
\[
\begin{aligned}
\bm{\varepsilon}^{(b)}&:=\nabla (\bm{\WarpShearIesan}\cdot\bm{b}_{\Section})
-\nu\big(\operatorname{dev}(\bm{r}\otimes\bm{r})-\bm{r}\otimes\CentroidLocation\big)\bm{b}_{\Section}, \\
\end{aligned}
\]
where \(\bm{\varepsilon}_{\vartheta}\) and \(\bm{\varepsilon}^{(b)}\) are the in–section vectors defined above. 
Equivalently, one may write the isotropic part using \((E,\nu)\) via
\(\displaystyle \bm{\sigma}=\frac{E}{1+\nu}\,\StrainTensor
+\frac{E\nu}{(1+\nu)(1-2\nu)}\,\operatorname{tr}(\StrainTensor)\,\mathbf{1}\).

\end{Quote}



Since $\FrameBasis\cdot\bm V_\Sigma=0$ (purely transverse) and $\operatorname{tr}\,\sym(\FrameBasis\otimes\bm V_\Sigma)=0$,
\[
\operatorname{tr}\StrainTensor
= \varepsilon\,\operatorname{tr}(\FrameBasis\otimes\FrameBasis)
- \nu\,\varepsilon\,\operatorname{tr}\oneS
= \varepsilon - 2\nu\,\varepsilon
= (1-2\nu)\,\varepsilon.
\]


From the constitutive equation \eqref{eq:linear-isotropy}
with  $\one=\FrameBasis\otimes\FrameBasis+\oneS$ and using $\operatorname{tr}\StrainTensor=(1-2\nu)\varepsilon$,
\[
\lambda(\operatorname{tr}\StrainTensor)\,\one
= \bigl(2G\nu\,\varepsilon\bigr)\,(\FrameBasis\otimes\FrameBasis+\oneS).
\]
Therefore
\[
\begin{aligned}
\StressTensor
&= 2G\Big[\varepsilon(\FrameBasis\otimes\FrameBasis)+\sym\{\FrameBasis\otimes \bm V_\Sigma\}-\nu\,\varepsilon\oneS\Big]
\;+\;2G\nu\,\varepsilon(\FrameBasis\otimes\FrameBasis+\oneS)\\[2pt]
&= 2G(1+\nu)\,\varepsilon\,(\FrameBasis\otimes\FrameBasis)\;+\;2G\,\sym\{\FrameBasis\otimes \bm V_\Sigma\}.
\end{aligned}
\]
Since $2G(1+\nu)=E$, an equivalent and compact form is
\[
\boxed{%
\;\;\StressTensor \;=\; E\,\varepsilon\,(\FrameBasis\otimes\FrameBasis)\;+\;2G\,\sym\{\FrameBasis\otimes \bm V_\Sigma\}\;.
}
\]
Fully expanded:
\[
\begin{aligned}
\StressTensor
&= E\Big[a_\varepsilon + (\bm{r}\cdot\IesanLoadM)
+ x\,\big((\bm{r}-\CentroidLocation)\cdot\bm{b}_{\Section}\big)\Big]\, \FrameBasis\otimes\FrameBasis
%
+\,G\Big(\FrameBasis\otimes \bm{\tau} + \bm{\tau} \otimes \FrameBasis\Big),
\end{aligned}
\]
with the transverse shear vector
\[
\bm{\tau} = G \bm V_\Sigma
% = (\FrameBasis\times\bm{r})\,(a_\vartheta+c_\vartheta)
% -\nu\bigl(\dev(\bm{r}\otimes\bm{r})-\bm{r}\otimes\CentroidLocation\bigr)\bm{b}_{\Section}
% + (\gradS\varphi)\,a_\vartheta
% +  c_\vartheta\,\gradS\WarpTwistIesan
% + 2\,\gradS\!\big(\bm{\WarpShearIesan}\,\bm{b}_{\Section}\big).
\]

% \begin{Quote}
% 
\textbf{1. Compute the Trace of the Strain Field}

We compute the trace of each term in the provided strain field $\bm{E}$.

\begin{enumerate}
    \item \textbf{Axial and Bending Terms ($E_A, E_B$):}
    \begin{align*}
    \operatorname{tr}\Big( (a_\varepsilon + \bm{r}\cdot\IesanLoadM)\,\big(\FrameBasis\otimes\FrameBasis-\nu\,\mathbf{1}_\Section\big) \Big)
    &= (a_\varepsilon + \bm{r}\cdot\IesanLoadM) \Big( \operatorname{tr}(\FrameBasis\otimes\FrameBasis) - \nu \operatorname{tr}\mathbf{1}_\Section \Big) \\
    &= (a_\varepsilon + \bm{r}\cdot\IesanLoadM)(1 - 2\nu)
    \end{align*}

    \item \textbf{Shear and Torsion Terms ($E_C, E_E$):}
    The vectors $\bm{\varepsilon}_{\vartheta}$ and $\bm{w}_b$ are defined on the cross-section and are thus orthogonal to $\FrameBasis$ (i.e., $\bm{\varepsilon}_{\vartheta} \cdot \FrameBasis = 0$ and $\bm{w}_b \cdot \FrameBasis = 0$).
    \[
    \operatorname{tr}\Big( \tfrac{1}{2}\big[\FrameBasis\otimes\bm{\varepsilon}_{\vartheta}+\bm{\varepsilon}_{\vartheta}\otimes\FrameBasis\big] \Big) = \tfrac{1}{2}(\FrameBasis \cdot \bm{\varepsilon}_{\vartheta} + \bm{\varepsilon}_{\vartheta} \cdot \FrameBasis) = 0
    \]
    \[
    \operatorname{tr}\Big( \tfrac{1}{2}\big[\FrameBasis\otimes\bm w_b+\bm w_b\otimes\FrameBasis\big] \Big) = \tfrac{1}{2}(\FrameBasis \cdot \bm{w}_b + \bm{w}_b \cdot \FrameBasis) = 0
    \]
    Thus, the shear terms are traceless.

    \item \textbf{Warping Normal Strain Term ($E_D$):}
    \begin{align*}
    \operatorname{tr}\Big( ((\bm{r}-\CentroidLocation)\cdot\bm{b}_{\Section})\,(\FrameBasis\otimes\FrameBasis) &\;-\;x\nu\,((\bm{r}-\CentroidLocation)\cdot\bm{b}_{\Section})\,\mathbf{1}_\Section \Big) \\
    &= ((\bm{r}-\CentroidLocation)\cdot\bm{b}_{\Section}) \operatorname{tr}(\FrameBasis\otimes\FrameBasis) - x\nu\,((\bm{r}-\CentroidLocation)\cdot\bm{b}_{\Section}) \operatorname{tr}(\mathbf{1}_\Section) \\
    &= ((\bm{r}-\CentroidLocation)\cdot\bm{b}_{\Section})(1) - x\nu\,((\bm{r}-\CentroidLocation)\cdot\bm{b}_{\Section})(2) \\
    &= ((\bm{r}-\CentroidLocation)\cdot\bm{b}_{\Section}) (1 - 2x\nu)
    \end{align*}
\end{enumerate}

Summing these components gives the total trace:
$$
\operatorname{tr}\StrainTensor = (1-2\nu)(a_\varepsilon + \bm{r}\cdot\IesanLoadM) + (1-2x\nu)((\bm{r}-\CentroidLocation)\cdot\bm{b}_{\Section})
$$

\textbf{2. Compute the Stress Field}

Separate the calculation into the normal stress components (proportional to $\FrameBasis \otimes \FrameBasis$ and $\mathbf{1}_\Section$) and the shear stress components.

\begin{itemize}
\item \textbf{Normal Stress Components}

Let $\bm{E}_{\text{normal}}$ be the non-shear terms from $\bm{E}$:
$$
\bm{E}_{\text{normal}} = (a_\varepsilon + \bm{r}\cdot\IesanLoadM)\,\big(\FrameBasis\otimes\FrameBasis-\nu\,\mathbf{1}_\Section\big) + ((\bm{r}-\CentroidLocation)\cdot\bm{b}_{\Section})\,(\FrameBasis\otimes\FrameBasis) \;-\;x\nu\,((\bm{r}-\CentroidLocation)\cdot\bm{b}_{\Section})\,\mathbf{1}_\Section
$$
The normal stress $\StressTensor_{\text{normal}}$ is:
$$
\StressTensor_{\text{normal}} = \lambda \operatorname{tr}(\bm{E}) \mathbf{1} + 2G \bm{E}_{\text{normal}}
$$
Group the terms multiplying $(\FrameBasis \otimes \FrameBasis)$ and $(\mathbf{1}_\Section)$:

\textbf{$(\FrameBasis \otimes \FrameBasis)$ terms:}
From $\lambda \operatorname{tr}(\bm{E}) \mathbf{1}$:
\begin{align*}
\lambda \operatorname{tr}(\bm{E}) (\FrameBasis \otimes \FrameBasis) &= \frac{E\nu}{(1+\nu)(1-2\nu)} \left( (1-2\nu)(a_\varepsilon + \bm{r}\cdot\IesanLoadM) + (1-2x\nu)((\bm{r}-\CentroidLocation)\cdot\bm{b}_{\Section}) \right) \FrameBasis \otimes \FrameBasis \\
&= \left( \frac{E\nu}{1+\nu} (a_\varepsilon + \bm{r}\cdot\IesanLoadM) + \frac{E\nu(1-2x\nu)}{(1+\nu)(1-2\nu)}((\bm{r}-\CentroidLocation)\cdot\bm{b}_{\Section}) \right) (\FrameBasis \otimes \FrameBasis)
\end{align*}
From $2G \bm{E}_{\text{normal}}$ using $2G = E/(1+\nu)$:
\begin{align*}
2G \left( (a_\varepsilon + \bm{r}\cdot\IesanLoadM) + ((\bm{r}-\CentroidLocation)\cdot\bm{b}_{\Section}) \right) (\FrameBasis \otimes \FrameBasis)
= \frac{E}{1+\nu} \left( (a_\varepsilon + \bm{r}\cdot\IesanLoadM) + ((\bm{r}-\CentroidLocation)\cdot\bm{b}_{\Section}) \right) (\FrameBasis \otimes \FrameBasis)
\end{align*}
Summing the $\FrameBasis \otimes \FrameBasis$ terms:
\begin{itemize}
    \item $(a_\varepsilon + \bm{r}\cdot\IesanLoadM)$: $\left( \frac{E\nu}{1+\nu} + \frac{E}{1+\nu} \right) = \frac{E(1+\nu)}{1+\nu} = E$
    \item $((\bm{r}-\CentroidLocation)\cdot\bm{b}_{\Section})$: $\left( \frac{E\nu(1-2x\nu)}{(1+\nu)(1-2\nu)} + \frac{E}{1+\nu} \right) = \frac{E}{1+\nu} \left( \frac{\nu(1-2x\nu) + (1-2\nu)}{1-2\nu} \right) = \frac{E}{1+\nu} \left( \frac{\nu-2x\nu^2 + 1-2\nu}{1-2\nu} \right) = \frac{E(1-\nu-2x\nu^2)}{(1+\nu)(1-2\nu)}$
\end{itemize}

\textbf{$\mathbf{1}_\Section$ terms:}
From $\lambda \operatorname{tr}(\bm{E}) \mathbf{1}$:
\[
\lambda (\operatorname{tr}\bm{E}) \mathbf{1}_\Section = \left( \frac{E\nu}{1+\nu} (a_\varepsilon + \bm{r}\cdot\IesanLoadM) + \frac{E\nu(1-2x\nu)}{(1+\nu)(1-2\nu)}((\bm{r}-\CentroidLocation)\cdot\bm{b}_{\Section}) \right) \mathbf{1}_\Section
\]
From $2G \bm{E}_{\text{normal}}$ using $2G = E/(1+\nu)$:
\begin{align*}
2G \left( -\nu(a_\varepsilon + \bm{r}\cdot\IesanLoadM) - x\nu((\bm{r}-\CentroidLocation)\cdot\bm{b}_{\Section}) \right) \mathbf{1}_\Section
= -\frac{E\nu}{1+\nu} \left( (a_\varepsilon + \bm{r}\cdot\IesanLoadM) + x((\bm{r}-\CentroidLocation)\cdot\bm{b}_{\Section}) \right) \mathbf{1}_\Section
\end{align*}
%
Summing the $\mathbf{1}_\Section$ terms:
% \begin{itemize}
%     \item $(a_\varepsilon + \bm{r}\cdot\IesanLoadM)$: $\left( \frac{E\nu}{1+\nu} - \frac{E\nu}{1+\nu} \right) = 0$
%     \item $((\bm{r}-\CentroidLocation)\cdot\bm{b}_{\Section})$: $\left( \frac{E\nu(1-2x\nu)}{(1+\nu)(1-2\nu)} - \frac{E\nu x}{1+\nu} \right) = \frac{E\nu}{1+\nu} \left( \frac{1-2x\nu - x(1-2\nu)}{1-2\nu} \right)$ \\
%     $= \frac{E\nu}{1+\nu} \left( \frac{1-2x\nu - x+2x\nu}{1-2\nu} \right) = \frac{E\nu(1-x)}{(1+\nu)(1-2\Vnu)}$
% \end{itemize}

\item \textbf{Shear Stress Components}

The shear terms in $\bm{E}$ are traceless, so the volumetric stress $\lambda \operatorname{tr}(\bm{E}) \mathbf{1}$ does not contribute to them. The shear stress is given solely by the $2G \bm{E}_{\text{shear}}$ term.
$$
\begin{aligned}
\StressTensor_{\text{shear}} &= 2G \bm{E}_{\text{shear}} \\
&= 2G \left( \tfrac12\left(\FrameBasis\otimes\bm{\varepsilon}_{\vartheta}+\bm{\varepsilon}_{\vartheta}\otimes\FrameBasis\right)(a_\vartheta+c_\vartheta) + \tfrac12\left(\FrameBasis\otimes\bm w_b+\bm w_b\otimes\FrameBasis\right) \right) \\
&= G(a_\vartheta+c_\vartheta)\left(\FrameBasis\otimes\bm{\varepsilon}_{\vartheta}+\bm{\varepsilon}_{\vartheta}\otimes\FrameBasis\right) + G\left(\FrameBasis\otimes\bm w_b+\bm w_b\otimes\FrameBasis\right)
\end{aligned}
$$
Using $\bm{\varepsilon}_{\vartheta}= \,\FrameBasis\times\bm{r} +\nabla \varphi$

\end{itemize}

\textbf{3. Final Stress Field}

Combining all components, the final stress tensor $\StressTensor$ is:
$$
\begin{aligned}
\StressTensor &=
\left( E (a_\varepsilon + \bm{r}\cdot\IesanLoadM) + (1-\nu-2x\nu^2)\frac{E}{(1+\nu)(1-2\nu)} ((\bm{r}-\CentroidLocation)\cdot\bm{b}_{\Section}) \right) \FrameBasis \otimes \FrameBasis \\[4pt]
&\quad + (1-x)\lambda \left( ((\bm{r}-\CentroidLocation)\cdot\bm{b}_{\Section}) \right) \mathbf{1}_\Section \\[4pt]
&\quad + G(a_\vartheta+c_\vartheta)\left(\FrameBasis\otimes\bm{\varepsilon}_{\vartheta}+\bm{\varepsilon}_{\vartheta}\otimes\FrameBasis\right) \\
&\quad + G\left(\FrameBasis\otimes\bm w_b+\bm w_b\otimes\FrameBasis\right)
\end{aligned}
$$


% \end{Quote}

\paragraph{Equilibrium}

Note that in the stress field from any superposition of Sain-Venant problems thus satisfies \eqref{eq:saint-venant-stress} and the equilibrium equation $\Div(\StressTensor) = \mathbf{0}$ yields two conditions:
\begin{subequations}
\begin{align}
    \ShearStress' &= \mathbf{0}  &\text{in }\Section \label{eq:shear_x} \\
    \DivS \ShearStress + \AxialStress' &= 0 &\text{in }\Section \label{eq:shear_div} \\
    \ShearStress \cdot \bm{\nu} &= 0 &\text{on } \SectionBoundary
\end{align}
\label{eq:sv-equilibrium}
\end{subequations}
From \eqref{eq:shear_x}, the shear stress $\ShearStress$ is independent of $\reflenx$, so $\ShearStress = \ShearStress(\bm{r})$. 
From \eqref{eq:shear_div}:
\begin{equation*}
    \DivS \ShearStress = -\AxialStress' 
    = -E ({\bm{b}_{\Section}} \cdot \bm{r} + b_{\varepsilon}) = -E(\bm{r}-\CentroidLocation)\cdot\bm{b}_{\Section}
\end{equation*}


\paragraph{Remarks.}

\begin{itemize}
\item The in-plane normal stress is identically zero:
$\StressTensor:\oneS=\mathbf{0}$, consistent with the imposed planar strain
$-\nu\,\varepsilon\oneS$.
\item Axial normal stress is purely uniaxial:
$\StressTensor:\FrameBasis\otimes\FrameBasis=E\,\varepsilon$
\end{itemize}

\begin{ambox}[Continuum Solution]{box:iesan}
\begin{equation*}
\begin{aligned}
  \hat{\bm{u}}\{\bm{a},\bm{b}_{\Section}\}
    &= 
      \left(-\tfrac{1}{2}x^2 \mathbf{1} - \nu\operatorname{dev}\bm{r}\otimes\bm{r} + x\, \FrameBasis\otimes\bm{r}\right)\IesanLoadM \\ 
    &+ (x\FrameBasis - \nu \bm{r})a_{\varepsilon} \\
    &+(x\; \FrameBasis\times \bm{r} + \varphi \FrameBasis)a_{\vartheta}\\
    & + (x\,\FrameBasis\times\bm{r}+ \WarpTwistIesan\FrameBasis)\, c_{\vartheta} \\
    &-\tfrac{1}{6}\,x^3\,\bm{b}_{\Section}
      - x\nu \, \bigl(\operatorname{dev}\bm{r}\otimes\bm{r} - \bm{r}\otimes\CentroidLocation\bigr)\;\bm{b}_{\Section} \\
    &+\tfrac{1}{2}x^2\,\FrameBasis\otimes\big(\bm{r} -\CentroidLocation\big)\bm{b}_{\Section}
    + \FrameBasis\otimes\bm{\WarpShearIesan}(\bm{r})\;\bm{b}_{\Section} .
\end{aligned}
% \label{eq:iesan-general-solution}
\end{equation*}

The infinitesimal strain is
\[
\StrainTensor
= \varepsilon\,\FrameBasis\otimes\FrameBasis
\;+\;\sym\{\FrameBasis\otimes \bm V_\Sigma\}
\;-\;\nu\,\varepsilon\,\oneS,
\]
where
\begin{equation*}
\begin{aligned}
\varepsilon
&= a_\varepsilon + \bm{r}\cdot\IesanLoadM
+ x\,\big((\bm{r}-\CentroidLocation)\cdot\bm{b}_{\Section}\big),
\\
\bm V_\Sigma
&= (\FrameBasis\times\bm{r} + \gradS \WarpTwistIesan)\,(a_\vartheta+c_\vartheta)
\;-\;\nu\bigl(\dev(\bm{r}\otimes\bm{r})-\bm{r}\otimes\CentroidLocation\bigr)\bm{b}_{\Section}
\;+\; \gradS \big(\bm{\WarpShearIesan}\cdot\bm{b}_{\Section}\big).
\end{aligned}
\end{equation*}

\[
\begin{aligned}
\StressTensor
&= E\varepsilon\, \FrameBasis\otimes\FrameBasis
%
+ 2\,G\Big(\FrameBasis\otimes \bm{\varepsilon}_{\gamma} + \bm{\varepsilon}_{\gamma} \otimes \FrameBasis\Big),
\end{aligned}
\]
\end{ambox}